
\section{（拟凝聚）代数层的相对 $\Spec$}

给定一个 $A$-代数 $R$，我们可以取它的 $\Spec$，从而得到 $\Spec A$ 上的一个\trAffineScheme：$\Spec R\to\Spec A$。现在我们将看到这种记号之整体化的一个\trUniversalProperty{}描述。考虑任意一个\trScheme{} $X$，及其上的一个\trQC{}\trSheafOfAlgebras（quasicoherent sheaf of algebras）$\mathscr{R}$。我们将定义如何取这个\trSheafOfAlgebras{}的 $\mathit{Spec}$，并且我们将得到一个「在 $X$ 上仿射」的\trScheme{} $\mathit{Spec}\mathscr{R}\to X$，也就是说，其\trStructureMorphism{}是一个\trAffine\trMorphism。你可以通过两种方式来领会它。

\textbf{17.1.1.}\quad 首先，也是最具体地：对任何\trAffine{}开集 $\Spec A\subset X$，$\Gamma(\Spec A,\mathscr{R})$ 是某个 $A$-代数；称其为 $R$。于是在 $\Spec A$ 之上，$\mathit{Spec}\mathscr{R}$ 将会是 $\Spec R$。

\textbf{17.1.2.}\quad 其次，它将满足一个\trUniversalProperty。我们本可以凭借如下事实来定义 $A$-\trScheme{} $\Spec R$：到 $\Spec R$ 的\trMorphism（从 $\Spec A$ 上的 $A$-\trScheme{} $W$ 出发）与 $A$-代数的映射 $R\to\Gamma(W,\OO_W)$ 相对应（这是我们的老朋友\mbox{练习~7.3.G}）。$\beta\colon\mathit{Spec}\mathscr{R}\to X$ 的\trUniversalProperty{}推广了这一点。给定一个\trMorphism{} $\mu\colon W\to X$，$X$-态射 $W\to\mathit{Spec}\mathscr{R}$ 与使交换图
\begin{equation*}\tag{17.1.2.1}\label{eq:relative-spec-univ-diagram}
\begin{tikzcd}[row sep=1.7em, column sep=2.6em]
& \OO_X \ar[ld] \ar[rd] & \\
\mathscr{R} \ar[rr, "\alpha"] && \mu_*\OO_W
\end{tikzcd}
\end{equation*}
得以交换的诸\trSheafOfAlgebras{}之间的\trMorphism{} $\alpha$ 以（对 $W$）函子的方式一一对应。这里，映射 $\OO_X\to\mu_*\OO_W$ 来自赋环空间映射，而映射 $\OO_X\to\mathscr{R}$ 则来自 $\mathscr{R}$ 的 $\OO_X$-代数结构。（致专家：$\mu_*\OO_W$ 未必是\trQC{}的，但这无关紧要。非专家请完全无视这条括注说明。）

凭一通\trUniversalProperty{}式的废话，这组数据便把 $\beta\colon\mathit{Spec}\mathscr{R}\to X$ 在相差唯一\trIsomorphism{}的意义下确定了下来——前提是它存在。

花哨的翻译：在 $X$-\trScheme{}的\trCategory{}中，$\beta\colon\mathit{Spec}\mathscr{R}\to X$ 表示函子
\[
(\mu\colon W\to X)\longmapsto\{(\alpha\colon\mathscr{R}\to\mu_*\OO_W)\}.
\]

\textbf{17.1.A.（练习）}\quad 证明：若 $X$ 是\trAffine{}的，比如说 $\Spec A$，而 $\mathscr{R}=\widetilde{R}$，其中 $R$ 是一个 $A$-代数，则 $\Spec R\to\Spec A$ 满足这个\trUniversalProperty。（提示：\mbox{练习~7.3.G}。）

\textbf{17.1.3.\ \trhProposition{}。}\ ——\ \emph{设 $\beta\colon\mathit{Spec}\mathscr{R}\to X$ 满足 $(X,\mathscr{R})$ 的\trUniversalProperty，而 $i\colon\mathfrak{U}\hookrightarrow X$ 是一个\trOpenSubscheme。那么 $\beta|_{\mathfrak{U}}\colon\mathit{Spec}\mathscr{R}\times_X\mathfrak{U}=(\mathit{Spec}\mathscr{R})|_{\mathfrak{U}}\to\mathfrak{U}$ 便满足 $(\mathfrak{U},\mathscr{R}|_{\mathfrak{U}})$ 的\trUniversalProperty。}

\emph{证明。}\quad 设 $W$ 是一个 $\mathfrak{U}$-\trScheme，带有\trStructureMorphism{} $\nu\colon W\to\mathfrak{U}$。$\mathfrak{U}$-态射 $W\to\mathit{Spec}\mathscr{R}\times_X\mathfrak{U}$ 与 $X$-态射 $W\to\mathit{Spec}\mathscr{R}$ 是同一回事（这里我们已经由假设知道 $\mu=i\circ\nu\colon W\to X$ 经由 $\mathfrak{U}$ 分解）：
\begin{tikzcd}[row sep=2.3em, column sep=2.6em]
W \ar[rd] \ar[ddr, bend right=28, "\nu"'] & & \\
& \mathit{Spec}\mathscr{R}\times_X\mathfrak{U} \ar[r] \ar[d] & \mathit{Spec}\mathscr{R} \ar[d, "\beta"] \\
& \mathfrak{U} \ar[r, hook, "i"] & X
\end{tikzcd}
由 $\mathit{Spec}\mathscr{R}$ 的\trUniversalProperty，这与一个映射 $\mathscr{R}\to\mu_*\OO_W$ 是同一信息；而由 $(\mu^*,\mu_*)$ 的\trAdjointness（\S14.5.5），它与 $\mu^*\mathscr{R}\to\OO_W$ 是同一回事；又（由 $\nu^*(\mathscr{R}|_{\mathfrak{U}})=\nu^*i^*\mathscr{R}=\mu^*\mathscr{R}$）它就是 $\nu^*(\mathscr{R}|_{\mathfrak{U}})\to\OO_W$。再由 $(\nu^*,\nu_*)$ 的\trAdjointness，这与 $\mathscr{R}|_{\mathfrak{U}}\to\nu_*\OO_W$ 是同一回事，\trUniversalProperty{}得证。\hfill$\square$

把\mbox{练习~17.1.A}与\mbox{命题~17.1.3}结合起来，我们便在如下情形证明了 $\mathit{Spec}\mathscr{R}$ 的存在性：$Y$ 是一个\trAffineScheme{}的\trOpenSubscheme，而 $\mathscr{R}$ 是那个\trAffineScheme{}上某个\trQC{}\trSheafOfAlgebras{}限制到 $Y$ 上的结果。

\textbf{17.1.B.（练习）}\quad 证明 $\mathit{Spec}\mathscr{R}$ 在一般情形下的存在性，其思路循着我们构造\trFiberProduct、\trNormalization{等等}的那套哲学。

\textbf{17.1.4.}\quad 我们来做几条快速的观察。首先，$\mathit{Spec}\mathscr{R}$ 可以「在 $X$ 上\trAffineLocal{}地计算」。我们还有一条\trIsomorphism{} $\phi\colon\mathscr{R}\overset{\sim}{\longrightarrow}\beta_*\OO_{\mathit{Spec}\mathscr{R}}$。

\textbf{17.1.C.（练习）}\quad 给定一个 $X$-态射 $\gamma$，
\begin{tikzcd}[row sep=2.2em, column sep=2.8em]
W \ar[rr, "\gamma"] \ar[dr, "\mu"'] && \mathit{Spec}\mathscr{R} \ar[dl, "\beta"] \\
& X, &
\end{tikzcd}
证明 \eqref{eq:relative-spec-univ-diagram} 中的那个态射 $\alpha$ 是合成
\[
\mathscr{R}\overset{\Phi}{\longrightarrow}\beta_*\OO_{\mathit{Spec}\mathscr{R}}\longrightarrow\beta_*\gamma_*\OO_W=\mu_*\OO_W.
\]

$\mathit{Spec}$ 构造给了我们一条理解\trAffine\trMorphism{}的重要途径。注意 $\mathit{Spec}\mathscr{R}\to X$ 是一个\trAffine\trMorphism。而「逆命题」也成立：

\textbf{17.1.D.（练习）}\quad 证明：若 $\mu\colon Z\to X$ 是一个\trAffine\trMorphism，则作为 $X$-\trScheme，我们有自然的\trIsomorphism{} $Z\overset{\sim}{\longleftarrow}\mathit{Spec}\mu_*\OO_Z$。

于是，借这个 $\mathit{Spec}$ 构造，我们可以复原任何\trAffine\trMorphism。更精确地说：一个\trMorphism{}是\trAffine{}的，当且仅当它具有上文所构造的 $\mu\colon\mathit{Spec}\mathscr{R}\to X$ 这种形式。

\textbf{17.1.E.（练习）}\quad 设 $\mu\colon\mathit{Spec}\mathscr{R}\to X$ 是这样一个\trMorphism。证明：$\mathit{Spec}\mathscr{R}$ 上诸\trQCSheaf{}的\trCategory{}等价于 $X$ 上带有 $\mathscr{R}$-\trModule{}结构的诸\trQCSheaf{}的\trCategory（即 $X$ 上拟凝聚的诸 $\mathscr{R}$-\trModule）。证明：$\mathit{Spec}\mathscr{R}$ 上有限型的诸\trQCSheaf{}之\trCategory{}等价于 $X$ 上带有「有限型」$\mathscr{R}$-\trModule{}结构的诸\trQCSheaf{}的\trCategory（即 $X$ 上有限型的拟凝聚 $\mathscr{R}$-\trModule）。

当 $X$ 十分简单而 $\mathit{Spec}\mathscr{R}$ 却颇为复杂时，这就很有用了。当 $X\cong\mathbb{P}^1$ 而 $\mathit{Spec}\mathscr{R}$ 是一条更复杂的\trCurve{}时，我们就会用到这一点；更一般地，在第 29 章证明\trSerreDuality{}时也会用到。

\textbf{17.1.F.（$\mathit{Spec}$ 的构造「与基变换交换」）}\quad 设 $\mu\colon Z\to X$ 是任意一个\trMorphism，而 $\mathscr{R}$ 是 $X$ 上一个\trQC{}\trSheafOfAlgebras。证明：存在自然的\trIsomorphism{} $Z\times_X\mathit{Spec}\mathscr{R}\overset{\sim}{\longleftarrow}\mathit{Spec}\mu^*\mathscr{R}$。（如果你琢磨一下，就会看到这推广了\mbox{命题~17.1.3} 的陈述，或许还有其证明。）

\textbf{17.1.5.}\ \emph{\trhDefinition{}。}\quad $\mathit{Spec}$ 构造的一个重要例子是\textbf{有限\trRank{}\trLocallyFreeSheaf{} $\mathscr{F}$ 的\trTotalSpace}（total space），我们把它定义为 $\mathit{Spec}(\operatorname{Sym}^\bullet\mathscr{F}^{\vee})$。

\textbf{17.1.G.（练习）}\quad 设 $\mathscr{F}$ 是\trRank{}为 $n$ 的一个\trLocallyFreeSheaf。证明：$\mathscr{F}$ 的\trTotalSpace{}是一个\trRank{}为 $n$ 的\emph{\trVectorBundle}，也就是说，给定任意一点 $p\in X$，都存在一个\trOpen\trNeighborhood{} $p\in\mathrm{U}\subset X$，使得
\[
\mathit{Spec}\bigl(\operatorname{Sym}^\bullet\mathscr{F}^{\vee}|_{\mathrm{U}}\bigr)\cong\mathbb{A}^n_{\mathrm{U}}.
\]
证明：$\mathscr{F}$ 同构于\trTotalSpace{} $\mathit{Spec}(\operatorname{Sym}^\bullet\mathscr{F}^{\vee})$ 的\trSection{}层。（可能的提示：使用\trTransitionFunction。）因此，\trTotalSpace{}也称作\textbf{与\trLocallyFreeSheaf{} $\mathscr{F}$ 相伴的\trVectorBundle}。（警示：一些作者，例如 [Stacks, tag 01M2]，把 $\mathit{Spec}(\operatorname{Sym}^\bullet\mathscr{F})$——即这个\trVectorBundle{}的\trDual——称作与 $\mathscr{F}$ 相伴的\trVectorBundle。）

特别地，若 $\mathscr{F}=\OO_X^{\oplus n}$，则 $\mathit{Spec}(\operatorname{Sym}^\bullet\mathscr{F}^{\vee})$ 就是 $\mathbb{A}^n_X$。

（旁注：你会注意到，构造 $\mathit{Spec}\operatorname{Sym}^\bullet$ 可以施用于任何\trCoherentSheaf{} $\mathscr{F}$（不取对偶，即 $\mathit{Spec}(\operatorname{Sym}^\bullet\mathscr{F})$）。这有时称作与 $\mathscr{F}$ 相伴的\textbf{\trAbelian\trCone}（abelian cone）。这个概念可能有用，但我们用不上它。）

\textbf{17.1.H.（$\mathbb{P}^n$ 上的\trTautologicalBundle{}是 $\OO(-1)$）}\quad 设 $k$ 是一个\trField。定义子集 $X\subset\mathbb{A}^{n+1}_k\times\mathbb{P}^n_k$ 为「$\mathbb{A}^{n+1}_k$ 中落在 $\mathbb{P}^n_k$ 的对应直线上的那些点」，使得映射 $\pi\colon X\to\mathbb{P}^n$ 在点 $\mathfrak{l}=[x_0,\cdots,x_n]$ 之上的\trFiber，恰是 $\mathbb{A}^{n+1}_k$ 中对应于 $\mathfrak{l}$ 的那条直线，也就是说，是 $(x_0,\dots,x_n)$ 的诸标量倍数。证明：$\pi\colon X\to\mathbb{P}^n_k$ 是（相应于\trInvertibleSheaf{} $\OO(-1)$ 的）那个\trLineBundle。（可能的提示：先在 $\mathbb{P}^n_k$ 的那些惯常的\trAffine{}开集上做，并弄清\trTransitionFunction。）因此，$\OO(-1)$ 常被称作 $\mathbb{P}^n_k$ 的\textbf{\trTautologicalBundle}（tautological bundle）（甚至在任意基上也是如此，而不只是域上）。（旁注：投射 $X\to\mathbb{A}^{n+1}_k$ 是 $\mathbb{A}^{n+1}_k$ 在「原点」处的\trBlowup；见\mbox{练习~10.3.F}。）

\textbf{17.1.6.\ 留待日后阅读：古怪的函子 $\pi^{!?}$，\trPushforward{}的一个偶现的右伴随。}\footnote{译注：上标「!?」整体是原书记号的一部分，并非排印笔误（页图放大核验）。}

由于接下来的论题会在两处不同的地方（\mbox{棘手练习~18.8.C} 与 \S29.3）被用到，而你或许会读、或许不会读它们，这段讨论除了放在这里，也再无其他自然的位置。但你现在不要去读它，而要等到需要它的时候再读，否则它将显得完全没有动机。

现在，当 $\pi$ 是一个\trFinite\trMorphism{}时，我们来为\trPushforward{} $\pi_*$ 定义一个\emph{右伴随} $\pi^{!?}$。（警示：记号 $\pi^{!?}$ 是非标准的，引入它只是出于我们将要给出的那些论证的需要。）这是令人惊讶的，因为我们通常把 $\pi_*$ 想成一个\emph{右伴随}（拉回 $\pi^*$ 的右伴随），而非左伴随。这可以粗略地类比于拉回那个出人意料的偶现\emph{左伴随}：\trExtensionByZero（\S23.4.7）。

\textbf{17.1.7.}\ \emph{\trhDefinition{}。}\quad 这个定义很简单，但很古怪。设 $\pi\colon X\to Y$ 是\trLocally\trNoetherian{}概形之间的一个\trFinite\trMorphism。若 $\mathscr{G}$ 是 $Y$ 上一个\trQCSheaf，则 $\mathscr{H}om_Y(\pi_*\OO_X,\mathscr{G})$ 也是 $Y$ 上一个\trQCSheaf（因为 $\pi_*\OO_X$ 是\trFinitelyPresented{}的，\mbox{练习~14.3.A(b)}）。它进一步带有 $(\pi_*\OO_X)$-\trModule{}的结构，这意味着它对应于 $X$ 上一个\trQCSheaf（由\mbox{练习~17.1.E}，该练习表明这个对应关系对 $\mathscr{G}$ 是函子性的）。令 $\pi^{!?}\mathscr{G}$ 为这个\trQCSheaf。那么 $\pi^{!?}$ 给出一个协变函子 $\mathit{QCoh}_Y\to\mathit{QCoh}_X$，并且也给出 $\mathit{Coh}_Y\to\mathit{Coh}_X$。

\textbf{17.1.I.（练习）}\quad 描述 $Y$ 上诸\trQCSheaf{}的一个自然\trIsomorphism
\[
\pi_*\pi^{!?}\mathscr{G}\overset{\sim}{\longrightarrow}\mathscr{H}om_Y(\pi_*\OO_X,\mathscr{G}),
\]
它对 $\mathscr{G}$ 是函子性的。等价地，描述函子的一个\trIsomorphism{} $\mathit{QCoh}_Y\to\mathit{QCoh}_Y$：
\[
\pi_*\pi^{!?}(\cdot)\overset{\sim}{\longrightarrow}\mathscr{H}om_Y(\pi_*\OO_X,\cdot).
\]

\textbf{17.1.8.}\quad 为了证明 $\pi^{!?}$ 是 $\pi_*$ 的右伴随，我们研究\trAffine{}情形，或者等价地，环论版本。设 $B\to A$ 是一个\trRingMorphism，$M$ 是一个 $A$-\trModule，$N$ 是一个 $B$-\trModule。注意 $\Hom_B(A,N)$ 自然带有 $A$-\trModule{}的结构。另外，$M$ 也自然带有 $B$-\trModule{}的结构；当我们想强调它作为 $B$-\trModule{}的结构时，有时便把它记作 $M_B$（见\mbox{练习~1.4.E}）。

考虑这个（十分奇怪的）映射
\begin{equation*}\tag{17.1.8.1}\label{eq:module-adjunction-map}
\Hom_A(M,\Hom_B(A,N))\longrightarrow\Hom_B(M_B,N)
\end{equation*}
其定义如下。给定 $m\in M$，以及 $\Hom_A(M,\Hom_B(A,N))$ 的一个元素 $\phi$，把 $m$ 送到 $\phi_m(1)$。

\textbf{17.1.J.（练习）}

（a）证明 \eqref{eq:module-adjunction-map} 是诸 $B$-\trModule{}之间的一个\trHomomorphism。

（b）证明 \eqref{eq:module-adjunction-map} 是一个双射。于是，$(M\mapsto M_B,\ N\mapsto\Hom_B(A,N))$ 构成一个\trAdjointPair{} $\mathit{Mod}_A\leftrightarrow\mathit{Mod}_B$。

\textbf{17.1.K.（练习）}\quad 设 $\pi\colon X\to Y$ 是\trLocally\trNoetherian{}概形之间的一个\trFinite\trMorphism。

（a）解释：前一练习中由 $N\mapsto\Hom_B(A,N)$ 给出的函子 $\mathit{Mod}_B\to\mathit{Mod}_A$ 如何就是\mbox{定义~17.1.7} 中 $\pi^{!?}$ 的\trAffineLocal{}描述。

（b）证明 $(\pi_*,\pi^{!?})$ 构成 $\mathit{QCoh}_X$ 与 $\mathit{QCoh}_Y$ 之间的一个\trAdjointPair。

（c）证明 $(\pi_*,\pi^{!?})$ 构成 $\mathit{Coh}_X$ 与 $\mathit{Coh}_Y$ 之间的一个\trAdjointPair。

（d）若 $\mathscr{F}\in\mathit{QCoh}_X$ 且 $\mathscr{G}\in\mathit{QCoh}_Y$，描述一个自然\trIsomorphism
\begin{equation*}\tag{17.1.8.2}\label{eq:sheaf-adjunction-isomorphism}
\pi_*\mathscr{H}om_X(\mathscr{F},\pi^{!?}\mathscr{G})\overset{\sim}{\longrightarrow}\mathscr{H}om_Y(\pi_*\mathscr{F},\mathscr{G})
\end{equation*}
（对 $\mathscr{F}$ 与 $\mathscr{G}$ 都是函子性的），它在\trAffineLocal{}上恰是\mbox{练习~17.1.J} 中所描述的那个\trIsomorphism{} \eqref{eq:module-adjunction-map}。
